% Section 7: Conclusion
% Paper: "Just a Glimpse: Temporal Wearing as a Design Paradigm for Transient Mixed Reality"

\section{Conclusion}

This paper has introduced \textit{temporal wearing} as a distinct interaction paradigm for mixed reality---one in which the act of donning and doffing a device constitutes the interaction itself, with encounters lasting seconds rather than sessions. We formalized this paradigm through a three-part wearing duration taxonomy that distinguishes temporal wear from permanent and session wear, and articulated a design space defined by body coupling, transition cost, and social legibility. Through six prototypes spanning handheld, neck-hung, and stationary form factors, we instantiated diverse positions within this design space, and through a field deployment at [VENUE PLACEHOLDER], we demonstrated that temporal wearing produces interaction patterns---double-takes, hand-offs, narration, and orbiting---that are qualitatively distinct from session-based MR use.

Our contributions are fourfold. First, the \textbf{temporal wearing concept and duration taxonomy} (C1) name a centuries-old but previously untheorized practice, identifying wearing duration as a primary design variable that existing wearable computing frameworks have overlooked. Second, the \textbf{design space} (C2) provides a structured vocabulary for reasoning about temporal wearing devices and situating them relative to existing MR technology. Third, our \textbf{empirical findings} (C3) establish that transition cost shapes re-engagement frequency, that temporal wearing enables socially shared MR without multi-user systems, and that visitors develop distinctive behavioral patterns around brief, repeated encounters. Fourth, our \textbf{six design principles} (C4)---Design for Discovery, Instant-On, Graceful Doff, Social Legibility, Invitation to Return, and Shared Glimpsing---offer actionable guidance for designers working with transient MR encounters.

More broadly, temporal wearing challenges the implicit assumption that longer wearing is better wearing. The history of optical instruments---lorgnettes, stereoscopes, viewfinders, View-Masters---demonstrates that brief, purposeful encounters with augmented perception have been valued for centuries. The contemporary pursuit of all-day, always-on wearable computing is a valid design direction, but it is not the only one. The full spectrum of wearing durations---from seconds to days---deserves design attention, and the brief end of that spectrum has been neglected.

We envision a future in which wearable computing embraces this full spectrum. Alongside devices designed to be worn and forgotten, there should be devices designed to be picked up and put down---devices that invite glimpses, reward returns, and support the social sharing of augmented perception. A visitor raises a viewer to her eyes. For a few seconds, the world is more than it appears. She lowers it, turns to her companion, and says: ``You have to see this.'' That moment---brief, shared, complete---is what temporal wearing is for.
